\scanpage{234}
\mspage{108}
% archive sheet number: 233

Le cas $b$ du corollaire est justiciable de la condition $c)$ du $1^\circ)$. (De plus $A$ ayant un objet final, est $W$-asph., donc on est dans les conditions du $3^\circ)$
Les cas $\deleted{\unsure{\mathord{\tikz[baseline=-.05ex,x=1em,y=1em]{\draw[line width=.055em] (.08,0)--(.78,.02)--(.69,.61)--(.02,.56)--cycle; \draw[line width=.05em] (.17,0)--(.11,.56);}}^\sim}}\mathord{\tikz[baseline=-.05ex,x=1em,y=1em]{\draw[line width=.055em] (0,0)--(.83,.03)--(.52,.67)--cycle; \draw[line width=.05em] (.13,.03)--(.45,.57);}},\mathord{\tikz[baseline=-.05ex,x=1em,y=1em]{\draw[line width=.055em] (0,0)--(.83,.03)--(.52,.67)--cycle; \draw[line width=.05em] (.13,.03)--(.45,.57);}}^{\sim},\mathord{\tikz[baseline=-.05ex,x=1em,y=1em]{\draw[line width=.055em] (.08,0)--(.78,.02)--(.69,.61)--(.02,.56)--cycle; \draw[line width=.05em] (.17,0)--(.11,.56);}}$ et $\mathord{\tikz[baseline=-.05ex,x=1em,y=1em]{\draw[line width=.055em] (.08,0)--(.78,.02)--(.69,.61)--(.02,.56)--cycle; \draw[line width=.05em] (.17,0)--(.11,.56);}}^\sim$, qui tombent en fait dans le cas $b)$ sont traités dans P.S. \S\,34 (p. 47--50)

\uline{Démonstration du th. 8.4.} La partie essentielle, techniquement, est bien sûr $1^\circ)$.$^\ast$ On a les implications $\pm$ tautologiques

% slanted
\margin{La partie $2^\circ$ en résulte, grâce à cor. 8.3.}

\[
 a\Longrightarrow b\Longrightarrow c\quad\deleted{\text{\unsure{immédiat|immédiate}}},
\]

où $a\Longrightarrow b$ provient du fait que $L$ est évidemment un objet injectif de $A^\wedge$. (L'objet de Lawvere d'un topos quelconque est un objet injectif). Il reste à établir

\[
 c\Longrightarrow d\Longrightarrow a.
\]

Que l'on ait $c\Longrightarrow d$ (et $c\Longrightarrow b$) est établi dans PS, voir le résumé des \unsure{critères|catég.} $W$-test de PS p. \unsure{168|468}--179\unsure{$'$} ; l'énoncé pertinent ici est prop. 3 page 173.

Il reste à établir $d)\Longrightarrow a)$ -- mais à présent on joue sur du velours !
Je \deleted{\illegible{1}} note ici le lemme suivant \deleted{\unsure{dont|qu'est}} la place serait plutôt à la fin du \S\,4 \unsure{des} Généralités sur les triples de Quillen :
